\scanpage{83}
\mspage{\handwritten{74}}
% One long diagonal cancellation crosses the main text of the sheet.
\deleted{ly constant with finite fibers and that the ${}_nG^{\unsure{o|0|\circ}}$ (for $n\handwritten{\geq}1$) are group schemes \deleted{\unsure{the}} finite locally free over $S$, that $H$ is finite locally free over $S$, and that \handwritten{t}he formation \deleted{commutes} of $\underline{\operatorname{Hom}}(\underline{\unsure{\varnothing|\phi|\Phi}}\ ,\ G^{\handwritten{o\natural}})$ commutes with\handwritten{\textsubscript{q}} any base change. Applying the same reflections to the extensions of $\underline{\varnothing}_n$ by $G^o_n$ , one finds that these correspond to the torsors under the Group $H_n=Hx_SS_n\handwritten{\simeq}\underline{\operatorname{Hom}}(\underline{\varnothing}_n,\ G^o_n)$. Now, $H$ being finite locally free over $S$, one immediately concludes from this that the \deleted{datum of the} category of torsors \inserted{$P$} under $H$ is equivalent (by the functor ``formal completion'') to the category of torsors under $\overset{\unsure{\scriptscriptstyle\wedge}}{H}$ i.e. of``coherent systems'' of torsors $P_n$ under the $H_n$ . This therefore implies that the category of extensions of $\underline{\varnothing}$ by $G^{o\handwritten{\natural}}$ is equivalent to that of the \deleted{extensions of} ``coherent systems'' of extensions of the $\underline{\varnothing}_n$ by the $G^o_n$ , which thus makes precise the definition of the extension $G^{\handwritten{\natural}}$ of (7.6.1).}

\deleted{The group scheme $G^{\handwritten{\natural}}$ , as an extension of $\underline{\varnothing}$ by $G^{o\handwritten{\natural}}$ , obviously still de\deleted{\uline{p\unsure{e}}}pends functorially on $G$. Its formation is compatible with base changes by a morphism of complete traits $S'\handwritten{\longrightarrow}S$.}

\margin{\handwritten{$\unsure{\checkmark|V}$}}
\deleted{7.6.3.Note that the \unsure{p|d}\unsure{r|e}eceding constructions also retain a meaning over a more general base than a complete \deleted{trait}\inserted{\uline{local} scheme}; it suffices to start with a noetherian adic ring $\deleted{\unsure{\mathbb{X}|\mathbb{V}|X}}\handwritten{V}$, whose topology is defined by an \deleted{sche} ideal $\underline{m}\handwritten{\subset}\deleted{\mathbb{X}}\handwritten{V}$, and a group scheme $G$ over $\deleted{\mathbb{X}}\handwritten{V}$ such that $G_{\unsure{o|0}}=G\otimes_{\deleted{\mathbb{X}}\handwritten{V}}\deleted{\mathbb{X}}\handwritten{V}_o$ (where $\deleted{\mathbb{X}}\handwritten{V}_o=\mathbb{X}/\underline{m}$) is an extension of a finite étale group scheme by an abelian scheme by a \deleted{scheme} torus. The essential question is that of showing that the abelian formal scheme $\overset{\handwritten{\scriptscriptstyle\wedge}}{B}$\inserted{over $\deleted{\mathbb{X}}\handwritten{V}$} constructed as in 7.1 in terms of $G^o$ is algebraizable; this is seen by the method of 7.2 when one assumes that $S$ is normal.}
% archive sheet number: 82
